\section{Cronograma de ejecuci\'on}

El cronograma organiza temporalmente las actividades del proyecto, garantizando coherencia con la metodolog\'ia CRISP-DM descrita en el Cap\'itulo II y con los objetivos planteados en el Cap\'itulo I.

\subsection{Per\'iodo}

El proyecto se ejecutar\'a en un per\'iodo de seis meses, comprendido entre marzo y agosto de 2025, con una dedicaci\'on promedio de 20 horas semanales por cada autor. Este horizonte permite completar la recolecci\'on de datos de Aduanet e INEI, el entrenamiento y comparaci\'on de modelos predictivos, el desarrollo del backend FastAPI y la aplicaci\'on React, la evaluaci\'on con productores de La Libertad, la redacci\'on del informe y la sustentaci\'on final.

\subsection{Cronograma}

La Tabla~\ref{tab:cronograma_gantt} presenta el cronograma de actividades en formato de diagrama de Gantt simplificado. Las marcas indican los meses en que se desarrolla cada actividad principal del proyecto.

\begin{table}[H]
\centering
\caption{Cronograma de ejecuci\'on del proyecto (diagrama de Gantt)}
\label{tab:cronograma_gantt}
\small
\begin{tabularx}{\textwidth}{p{5.5cm}*{6}{>{\centering\arraybackslash}X}}
\toprule
\textbf{Actividad} & \textbf{Mar} & \textbf{Abr} & \textbf{May} & \textbf{Jun} & \textbf{Jul} & \textbf{Ago} \\
\midrule
Validaci\'on de instrumentos de recolecci\'on & X & & & & & \\
Recolecci\'on de datos (Aduanet, INEI, encuestas) & X & X & & & & \\
An\'alisis y tratamiento de datos (ETL, limpieza) & & X & X & & & \\
Modelado predictivo (SARIMA, LSTM, SVR) & & X & X & & & \\
Desarrollo backend FastAPI y app React & & & X & X & & \\
Evaluaci\'on e interpretaci\'on de resultados & & & & X & X & \\
Redacci\'on del informe final & & & & X & X & X \\
Revisi\'on y correcciones & & & & & X & X \\
Sustentaci\'on & & & & & & X \\
\bottomrule
\end{tabularx}
\parbox{0.94\textwidth}{\small \textit{Nota.} Elaboraci\'on propia. X = actividad programada en el mes indicado. Per\'iodo: marzo--agosto 2025.}
\end{table}


Las actividades se articulan de la siguiente manera: en los meses 1 y 2 se validan instrumentos y se recolectan datos; en los meses 2 y 3 se realiza el an\'alisis, tratamiento y modelado; en los meses 3 y 4 se desarrolla la soluci\'on tecnol\'ogica; en los meses 4 y 5 se interpretan resultados y se eval\'ua con productores; y en los meses 5 y 6 se redacta, revisa y sustenta el informe final.
